% bridge-handout.tex
\documentclass[11pt, a4paper]{article}

\usepackage[T1]{fontenc}
\usepackage[utf8]{inputenc}
\usepackage{tgpagella}
\usepackage[spanish, es-noshorthands]{babel}
\usepackage{amsmath, amssymb, bm}
\usepackage[margin=2.5cm]{geometry}
\usepackage{fancyhdr}
\usepackage[most]{tcolorbox}

\pagestyle{fancy}
\fancyhf{}
\setlength{\headheight}{16pt}
\lhead{\textbf{UCB Sede Tarija} | Ing. Mecatrónica}
\rhead{Solucionario y Apoyo EC3: Puente Matemático[cite: 4]}
\lfoot{Prof: M.Sc. Ing. Bernardo Quiroga Turdera[cite: 1, 3]}
\renewcommand{\headrulewidth}{0.4pt}
\definecolor{ucbblue}{RGB}{0, 51, 102}

\begin{document}

\begin{tcolorbox}[colback=ucbblue!5!white, colframe=ucbblue, title=\textbf{Artefacto 2: Mini-Puente Matemático para Circuitos}]
    Este documento asiste en la interpretación física de la derivada antes de resolver EDOs formales[cite: 4].
\end{tcolorbox}

\section*{Sección A: Tasa de Cambio[cite: 4]}
Complete el significado físico:
\begin{itemize}
    \item $\frac{dx}{dt} > 0$ significa que la variable está \underline{\hspace{4cm}}.[cite: 4]
    \item $\frac{dx}{dt} = 0$ significa que la variable es \underline{\hspace{4cm}}.[cite: 4]
    \item $\frac{dx}{dt} < 0$ significa que la variable está \underline{\hspace{4cm}}.[cite: 4]
\end{itemize}

\section*{Sección B: Cambio Promedio[cite: 4]}
Dado que $x(1\,\text{s}) = 2$ y $x(3\,\text{s}) = 8$[cite: 4].\\
Encuentre la tasa promedio $\frac{\Delta x}{\Delta t}$: \underline{\hspace{4cm}}[cite: 4].

\section*{Sección C: Traducción Eléctrica[cite: 4]}
Interprete en palabras el significado de los siguientes valores instantáneos:
\begin{itemize}
    \item $\frac{dv_C}{dt} = +5\,\text{V/s}$: \underline{\hspace{8cm}}[cite: 4].
    \item $\frac{dv_C}{dt} = 0$: \underline{\hspace{8cm}}[cite: 4].
    \item $\frac{di_L}{dt} = -2\,\text{A/s}$: \underline{\hspace{8cm}}[cite: 4].
\end{itemize}

\section*{Sección D: Significado de la EDO[cite: 4]}
Dada la ecuación diferencial generada por KVL en un circuito RC:
$$RC \frac{dv_C}{dt} + v_C = V_s$$[cite: 4]

Identifique y etiquete:
\begin{itemize}
    \item La variable de estado actual: \underline{\hspace{4cm}}[cite: 4].
    \item El término que representa la tasa de cambio: \underline{\hspace{4cm}}[cite: 4].
    \item La excitación (entrada) del sistema: \underline{\hspace{4cm}}[cite: 4].
\end{itemize}

\end{document}
